\section{Erd\H{o}s Problem \#229: derivative zeros on prescribed discrete sets}
\noindent Partial reconstruction with a complete finite-set case, 23 September 2026.

\subsection*{1) Statement and result proved here}
Let $(S_n)_{n\ge1}$ be subsets of $\mathbb C$, each with no finite accumulation point. Is there a transcendental entire $f$ such that for each $n$ some derivative order $k_n\ge0$ satisfies $f^{(k_n)}|_{S_n}=0$? The full affirmative theorem is due to Barth and Schneider.

We prove the assertion when every $S_n$ is finite, allowing their union to be dense or to have accumulation points. We also treat the separate case where the whole union is locally finite. Neither condition covers the full question.

\subsection*{2) Polynomial corrections for finite sets}
Assume each $S_n$ is finite. Start with $F_0=0$; for this starting step only, use the degree bound $D_0=0$. At stage $n$, put $D_{n-1}=\deg F_{n-1}$ if $n>1$ and choose
\[
 k_n=D_{n-1}+1,\qquad t_n=D_{n-1}+1,
 \qquad W_n(z)=\prod_{i=1}^n\prod_{s\in S_i}(z-s)^{k_i+1}.
\]
Empty products are one. Choose a nonzero complex $\epsilon_n$ sufficiently small that
\[
 h_n(z)=\epsilon_n z^{t_n}W_n(z),\qquad
 \max_{|z|\le n}|h_n(z)|\le2^{-n},\tag{1}
\]
and set $F_n=F_{n-1}+h_n$. Such a choice is possible because $z^{t_n}W_n$ is a fixed nonzero polynomial with finite norm on the disk.

At each $s\in S_n$, the derivative $F_{n-1}^{(k_n)}$ vanishes identically, since $k_n>D_{n-1}$. The correction $h_n$ has a zero of order at least $k_n+1$ at $s$, so $h_n^{(k_n)}(s)=0$ too. Every later correction $h_m$, $m>n$, contains the factor $(z-s)^{k_n+1}$, since $W_m$ includes all earlier factors. Thus
\[
 F_m^{(k_n)}(s)=0\qquad(m\ge n,\ s\in S_n).\tag{2}
\]
Repeated occurrences of a point in different sets cause no difficulty: the multiplicities in the product only increase.

\subsection*{3) Entire limit and transcendence}
Estimate (1) makes $\sum h_n$ uniformly convergent on every fixed disk: after finitely many terms its norms there are bounded by the summable sequence $2^{-n}$. The limit $f$ is entire, and derivatives converge locally uniformly. Passing to the limit in (2) gives the requested vanishing for every $n$ in the finite-set case.

It remains to ensure that $f$ is not a polynomial, rather than merely nonzero. Let
\[
 d_n=t_n+\operatorname{ord}_0W_n
\]
be the lowest degree of a nonzero coefficient of $h_n$. Then $d_n>D_{n-1}$, so its coefficient in $F_n$ is nonzero and has no contribution from earlier stages. Every later correction has lowest degree strictly greater than the degree of the preceding partial polynomial, hence greater than or equal to $\deg F_n+1>d_n$. Consequently the coefficient of $z^{d_n}$ remains permanently nonzero in $f$. The strictly increasing sequence $(d_n)$ supplies infinitely many nonzero Taylor coefficients. This proves transcendence, including when zero belongs to some or all of the $S_n$.

\subsection*{4) A locally finite union, with arbitrary set sizes}
If $S=\bigcup_n S_n$ itself has no finite accumulation point, the Weierstrass product theorem supplies an entire function whose zero set is $S$. If $S$ is infinite, this nonzero function is automatically transcendental, since a nonzero polynomial has only finitely many zeros. Taking $k_n=0$ for every $n$ solves this case. If $S$ is finite, use $e^z\prod_{s\in S}(z-s)$; its zero set is $S$ and it is transcendental. For $S=\varnothing$, use $e^z$. This paragraph explicitly uses the classical Weierstrass product theorem as an external analytic input.

\subsection*{5) Remaining gap and sources}
\textbf{UNRESOLVED in this attempt.} The unhandled case permits infinite individual discrete sets whose union has finite accumulation points. The finite-set construction relies on every partial sum being a polynomial, so that a sufficiently high derivative annihilates it. Replacing $W_n$ by an infinite entire product destroys that property; choosing $k_n$ larger than a nonexistent finite degree is not a valid extension. Nor can one simply prescribe zeros of $f$ on a union with an accumulation point, by the identity theorem.

Source: \url{https://www.erdosproblems.com/229}, checked 23 September 2026, including its attribution of the full result to Barth and Schneider (1972). The finite-set construction is proved above. No reconstruction of the full infinite-set theorem or independent Lean verification is claimed.

The estimate measures coverage of the existence question, not a correctness probability.
\subsection*{6) COMPLETION ESTIMATE}
\textbf{COMPLETION: 50\%}
